% ============================================================
% Chapter 1: Introduction
% ============================================================
\section{Introduction}
\label{sec:introduction}

\subsection{Background and Motivation}

The framing of immigration in news media profoundly shapes public attitudes toward migrants, asylum seekers, and integration policies. Within Germany, the ongoing discourse surrounding immigration---intensified by the 2015--2016 European refugee crisis and subsequent waves from Ukraine in 2022---has produced a vast corpus of news coverage that embeds economic arguments both for and against immigration. Understanding how German media frames immigration as either an \textit{economic threat} or an \textit{economic benefit} is essential for political communication research, as these frames have been shown to influence public opinion, voting behaviour, and policy outcomes~[1, 6].

This project addresses the challenge of annotating economic framing at scale. Manual content analysis of 10,000 paragraphs would require hundreds of person-hours, introducing significant inter-coder reliability concerns. We adopt a computational inference approach using Large Language Models (LLMs), specifically the \texttt{Mistral-3-14B} model hosted on the University of Mannheim's GPU cluster, guided by the \textbf{bacchuss annotation routine} developed by Freudenthaler~[4].

\subsection{Research Question}

The central research question guiding this project is:

\begin{quote}
\textit{To what extent do German immigration news articles (2022--2026) frame immigration in terms of economic threat versus economic benefit, and how do these framing patterns vary across news outlets?}
\end{quote}

Each paragraph is evaluated along two separate binary dimensions:
\begin{itemize}
    \item \textbf{Economic Threat (T1--T5):} Does the paragraph frame immigration as causing economic harm, labour market disruption, welfare strain, or capacity overload?
    \item \textbf{Economic Benefit (B1--B5):} Does the paragraph frame immigration as generating economic value, filling labour shortages, or contributing to demographic sustainability?
\end{itemize}

Both dimensions are coded separately per paragraph, allowing detection of \textit{mixed framing}---paragraphs that invoke both threat and benefit arguments simultaneously.

\subsection{Project Scope and Objectives}

The objectives of this project were:

\begin{enumerate}
    \item \textbf{Dataset Preparation:} Draw a simple random sample of 10,000 paragraphs from a master corpus of 659,895 German news paragraphs, translate them from German to English for processing, and prepare a 200-row development sample for iterative validation.
    
    \item \textbf{Instruction Development:} Design and iteratively refine annotation instructions following the Freudenthaler bacchuss routine, incorporating sub-criteria from de~Vreese et~al.\ (2010)~[2] and Guo et~al.\ (2023)~[3].
    
    \item \textbf{Validation:} The prompt-development process targeted an F1-score of 0.80 for both dimensions, while allowing a documented exception before production. Threat passed the nominal pilot gate and Benefit was accepted as a documented exception at F1 0.769, after \textbf{8 full iterations} of prompt refinement against the pilot gold standard.
    
    \item \textbf{Production Annotation:} Annotate all 10,000 paragraphs with the selected final instruction set, producing structured explanations for each classification.
    
    \item \textbf{Analysis:} Generate frame distributions overall and by publication, identifying systematic framing patterns across the German media landscape.
\end{enumerate}

\subsection{Key Contributions}

This project makes several contributions to the methodology of computational content analysis:

\begin{itemize}
    \item A fully documented, reproducible annotation pipeline with 8 iterations of instruction refinement, each measured against a 200-row human gold standard.
    \item Detailed documentation of every failure, retry, and correction made during the development process, providing a transparent account of the iterative nature of LLM-assisted annotation.
    \item A production-scale annotated corpus of 10,000 paragraphs with structured evidence-based explanations.
    \item Cross-outlet comparison revealing significant framing asymmetries between right-leaning outlets (Jouwatch, Tichys Einblick) and business-oriented outlets (Handelsblatt).
\end{itemize}

\subsection{Report Structure}

This report is organised as follows. Section~\ref{sec:literature} reviews the theoretical foundations of framing analysis and computational annotation. Section~\ref{sec:methodology} describes the dataset, the bacchuss routine, and the operationalisation of economic framing dimensions. Section~\ref{sec:pipeline} details the complete technical pipeline from data preparation through production annotation. Section~\ref{sec:iterations} presents the full 8-iteration prompt development log, documenting every failure and retry in detail. Section~\ref{sec:results} reports the final corpus distribution and per-publication analysis. Section~\ref{sec:discussion} discusses the findings, limitations, and implications. Section~\ref{sec:conclusion} concludes with recommendations for future work.

\newpage
